\documentclass[11pt]{article}
\usepackage[T1]{fontenc}
\usepackage[utf8]{inputenc}
\usepackage[spanish,es-nodecimaldot,es-noshorthands]{babel}
\usepackage{lmodern}
\usepackage[a4paper,margin=24mm,headheight=15pt]{geometry}
\usepackage{amsmath,amssymb,booktabs,longtable,tabularx,array}
\usepackage{graphicx,xcolor,enumitem,fancyhdr,tikz,pgfplots}
\pgfplotsset{compat=1.18}
\usetikzlibrary{arrows.meta,positioning,babel}
\usepackage{csquotes}
\usepackage[backend=biber,style=apa,sorting=nyt]{biblatex}
\addbibresource{references.bib}
\definecolor{navy}{HTML}{17324D}
\definecolor{teal}{HTML}{007F82}
\definecolor{light}{HTML}{EDF3F6}
\usepackage[hypertexnames=false,colorlinks=true,linkcolor=navy,citecolor=teal,urlcolor=teal]{hyperref}
\usepackage{bookmark}
\setlength{\parindent}{0pt}
\setlength{\parskip}{6pt}
\setlength{\emergencystretch}{2em}
\setlist{nosep,leftmargin=*}
\pagestyle{fancy}\fancyhf{}
\fancyhead[L]{\small\textcolor{navy}{Robótica y Sistemas Autónomos}}
\fancyhead[R]{\small Borrador v0.2}
\fancyfoot[L]{\small 2026-2 | Revisión de Jairzinho}
\fancyfoot[R]{\thepage}
\newcommand{\pendiente}[1]{\par\noindent\textcolor{teal}{\textbf{Pendiente:} #1}\par}
\newcommand{\nota}[1]{\par\noindent\colorbox{light}{\parbox{\dimexpr\linewidth-2\fboxsep}{#1}}\par}
\newcommand{\portada}[2]{\begin{titlepage}\vspace*{17mm}{\Large\textcolor{teal}{TRABAJO PARCIAL · 2026-2}}\par\vspace{12mm}{\Huge\bfseries\color{navy}#1\par}\vspace{9mm}{\Large #2\par}\vfill{\large Robot antropomórfico de 3 GDL\par Robótica y Sistemas Autónomos\par}\vspace{10mm}Docente: Ing. Wilson Dennis Reyes Vásquez\par Fecha de preparación: 28 de octubre de 2026\par\vspace{8mm}\nota{Documento de trabajo para revisión. Los supuestos, pendientes y resultados preliminares no constituyen una entrega aprobada.}\end{titlepage}}

\begin{document}
\begin{titlepage}\centering
\vspace*{7mm}
\IfFileExists{../../assets/branding/uni-logo.png}{\includegraphics[width=30mm]{../../assets/branding/uni-logo.png}}{\IfFileExists{../../assets/branding/uni-logo.pdf}{\includegraphics[width=30mm]{../../assets/branding/uni-logo.pdf}}{\fbox{\parbox[c][24mm][c]{45mm}{\centering Logo oficial UNI\\pendiente}}}}
\par\vspace{6mm}{\Large UNIVERSIDAD NACIONAL DE INGENIERÍA\par}
\vspace{5mm}{\large [Nombre oficial de la maestría]\par}
\vspace{12mm}{\Huge\bfseries\color{navy}Trabajo Parcial\par}
\vspace{5mm}{\Large Modelado y simulación cinemática de un robot antropomórfico de 3 GDL\par}
\vspace{5mm}Basado en el robot VSL de Ramzy et al. (2024)\\Alcance rígido nominal pendiente de confirmación
\vfill
\begin{tabular}{ll}Curso:&Robótica y Sistemas Autónomos\\Docente:&Ing. Wilson Dennis Reyes Vásquez\\Semestre:&2026-2\\Autores:&Josemanuel [apellidos]\\&Melissa [apellidos]\\&Jairzinho [apellidos]\\&Renzo [apellidos]\\&Kenny [apellidos]\end{tabular}
\vfill Lima, Perú\par[Fecha de entrega]\par\vspace{6mm}
\nota{Plantilla v0.2 para revisión. El borrador de cinemática diferencial se mantiene separado. Completar datos de carátula y eliminar las notas de trabajo sólo después de revisión.}
\end{titlepage}
\tableofcontents\newpage
\section{Introducción}
\pendiente{Presentar el problema, importancia de la cinemática de un robot antropomórfico y propósito de la simulación. Sustentar el contexto con citas.}
\textbf{Objetivo general propuesto:} modelar y simular en Python la cinemática de un manipulador antropomórfico de 3 GDL basado en un robot documentado, integrando análisis directo, inverso, diferencial, singularidades y control cinemático.
\textbf{Objetivos específicos:} establecer DH y cinemática directa; obtener soluciones inversas; derivar y verificar Jacobianos; caracterizar singularidades de posición; evaluar una ley de control mediante métricas reproducibles.
\textbf{Alcance provisional:} geometría y cinemática rígida nominal del robot de \textcite{ramzy2024}. Confirmar si satisface la consigna \parencite{consigna2026}; no afirmar simulación de flexibilidad ni réplica experimental.

\section{Desarrollo}
\subsection{Robot seleccionado, hipótesis y convenciones}
\pendiente{Describir el prototipo del artículo, justificar su selección y presentar esquema propio con marcos. Documentar las correcciones de las matrices impresas.}
\begin{center}\begin{tabular}{lcl}\toprule Parámetro & Valor & Fuente\\\midrule
$l_1$ & 0.04 m & Paper, tabla 2\\ $l_2$ & 0.29 m & Paper, tabla 2\\ $l_3$ & 0.24 m & Paper, tabla 2\\\bottomrule\end{tabular}\end{center}
Contrato propuesto: DH estándar, base fija, radianes, metros, vectores columna y velocidad geométrica ordenada como $[v;\omega]$. Los límites físicos de articulaciones aún deben documentarse.

\subsection{Cinemática directa · Josemanuel}
\pendiente{Incluir tabla DH, matrices individuales válidas, producto $T(q)$, posición, orientación y comprobación de $R^TR=I$, $\det R=1$.}
\pendiente{Simular configuraciones representativas y espacio de trabajo en Python; distinguir rango de muestreo de límites físicos. Entregar figura y comparación entre DH y fórmula cerrada.}

\subsection{Cinemática inversa · Melissa}
\pendiente{Derivar usando altura $z-l_1$, funciones \texttt{atan2}, ramas de codo y condiciones de alcanzabilidad. Reconciliar la notación del paper antes de reutilizar ecuaciones.}
\pendiente{Evaluar objetivos alcanzables e inalcanzables y reconstrucción mediante directa. Reportar residuo cartesiano y tratamiento de singularidades.}

\subsection{Cinemática diferencial · Jairzinho}
\nota{Existe un borrador separado en \texttt{research/robot-3r/differential.md}. Su contenido aún no se incorpora a este informe.}
\pendiente{Tras revisión: incluir derivación de $J_v$ y $J_\omega$, convención de velocidad, ejemplo numérico y comparación con derivadas finitas en Python.}
\pendiente{Incorporar gráficas de velocidades, error y sensibilidad al paso temporal, con procedencia y unidades.}

\subsection{Singularidades · Renzo}
\pendiente{Definir la tarea de posición, analizar $\det J_v$, rango y valores singulares. Estudiar alineación de eslabones, eje de base e intersecciones.}
\pendiente{Distinguir singularidades del Jacobiano de posición de las del Jacobiano geométrico completo. Mostrar consecuencias sobre velocidades y control.}

\subsection{Control cinemático · Kenny}
\pendiente{Definir trayectoria deseada, condiciones iniciales, ley de control, ganancias, discretización, límites y regularización.}
\pendiente{Simular error de seguimiento, velocidades articulares y proximidad a singularidad; comparar al menos un escenario regular y uno exigente. No confundir movimiento prescrito con control cerrado.}

\subsection{Resultados integrados y discusión}
\begin{center}\begin{tabularx}{\linewidth}{@{}lXX@{}}\toprule
Bloque & Métrica sugerida & Evidencia pendiente\\\midrule
Directa & Residuo DH frente a fórmula & Figuras y registro Python\\
Inversa & $\|p(q)-p_d\|$ & Ramas y casos inalcanzables\\
Diferencial & $\|J_v\dot q-\dot p_{FD}\|$ & Curvas y convergencia temporal\\
Singularidades & $\sigma_{\min}$ y rango & Casos regulares y singulares\\
Control & RMSE, máximo error, $\max|\dot q_i|$ & Trayectoria y sensibilidad\\\bottomrule\end{tabularx}\end{center}
\pendiente{Presentar resultados ejecutados, método de cálculo, versiones y límites. Separar evidencia propia de datos del paper. Discutir la discrepancia de la tabla 3 sin atribuirle una causa no verificada.}

\section{Apreciaciones personales y conjunta}
\subsection{Josemanuel: cinemática directa}
\pendiente{Apreciación auténtica: aprendizaje, dificultad, decisión de modelado y límites.}
\subsection{Melissa: cinemática inversa}
\pendiente{Apreciación auténtica sobre ramas, alcanzabilidad y comprobación.}
\subsection{Jairzinho: cinemática diferencial}
\pendiente{Apreciación auténtica sobre Jacobianos, velocidades y validación.}
\subsection{Renzo: singularidades}
\pendiente{Apreciación auténtica sobre rango, tarea y consecuencias.}
\subsection{Kenny: control cinemático}
\pendiente{Apreciación auténtica sobre seguimiento, ganancias y límites.}
\subsection{Apreciación conjunta}
\pendiente{Síntesis consensuada de los cinco sobre integración y aprendizaje. Conservar ambas opciones hasta decidir el formato final.}

\section{Conclusiones}
\pendiente{Responder a cada objetivo usando resultados medidos. Distinguir consistencia matemática, desempeño de simulación y validación experimental. No escribir conclusiones de simulaciones aún no ejecutadas.}

\section{Recomendaciones}
\pendiente{Proponer mejoras concretas derivadas de los resultados: paso temporal, límites de operación, regularización, parámetros faltantes y posible extensión flexible. Justificar cada recomendación.}

\printbibliography[heading=bibnumbered,title={Referencias}]
\appendix
\section{Esquemas realizados}
\pendiente{Anexo obligatorio según la consigna: incorporar esquema propio del robot con articulaciones, longitudes, ejes y marcos DH. El borrador diferencial contiene un primer esquema para revisar.}
\section{Código y reproducibilidad}
\pendiente{Identificar scripts, orden de ejecución, versión de Python y bibliotecas empleadas, entradas, tolerancias y salidas. Referenciar archivos entregados, sin pegar código innecesario en el cuerpo.}
\section{Pruebas y observaciones}
\pendiente{Incluir matriz de pruebas, registros, incidencias del artículo y decisiones de alcance confirmadas por el profesor.}

\end{document}
